% ==============================================================================
% capitulos/05_implementacao_nodered.tex
% Capítulo 5 — Implementação da Pipeline Analítica no Node-RED e Persistência
% Inclui Diagrama TikZ 4, Equações Estatísticas, Quadros 2, 3, 4 e 5
% ==============================================================================

\chapter{Implementação da Pipeline Analítica no Node-RED e Persistência}
\label{cap:implementacao_nodered}

Este capítulo detalha a camada de Tecnologia da Informação (IT) responsável pela coleta contínua de dados via Modbus TCP/IP, armazenamento estruturado em PostgreSQL, análise estatística determinística em janelas temporais deslizantes, diagnóstico contextual por agente baseado em LLM com validação estrita por \textit{JSON Schema}, envio de comandos de atuação aos LEDs e emissão automática de relatórios operacionais.

\section{Arquitetura do Fluxo no Node-RED}
\label{sec:nodered_flow}

A plataforma \textbf{Node-RED} é executada localmente no servidor \texttt{192.168.10.20}, com os pacotes \texttt{node-red-contrib-modbus} e \texttt{node-red-contrib-postgresql} instalados via gerenciador de paleta (\textit{Manage Palette}) \cite{nodered_docs, nodered_contrib_modbus}. O nó de configuração compartilhado \texttt{Modbus-Client} é parametrizado para conexão TCP ao endereço \texttt{192.168.10.10}, porta \texttt{502}, \textit{Unit ID} \texttt{1}, com enfileiramento sequencial de comandos, \textit{timeout} de requisição de $1500\,\text{ms}$ e reconexão automática após $2000\,\text{ms}$ em caso de queda de \textit{socket}.

O fluxo no Node-RED é organizado em dois ramos cooperativos sincronizados pelo PostgreSQL e pelo contexto de aba (\texttt{flow}), conforme representado na Figura~\ref{fig:fluxo-pipeline-nodered}:
\begin{enumerate}
  \item \textbf{Ramo de Ingestão Rápida ($T_p = 500\,\text{ms}$):} os nós \texttt{Modbus-Read} (ou um nó \texttt{inject} de $500\,\text{ms}$ acionando nós \texttt{Modbus-Flex-Getter}) leem em bloco as tabelas \texttt{FC02} e \texttt{FC04}, um nó \texttt{function} estrutura o objeto \texttt{msg.payload} com carimbo de tempo UTC, grava o registro bruto na tabela \texttt{telemetria\_sensores} via nó \texttt{postgresql} e comanda um nó \texttt{Modbus-Flex-Write} (\texttt{FC06} em \texttt{HR[0]}) para incrementar o \textit{heartbeat} e manter o \textit{watchdog} do PLC500 ativo.
  \item \textbf{Ramo Analítico Híbrido e Atuação ($W = 60\,\text{s}$ ou disparo imediato por evento crítico):} consulta as amostras da última janela temporal no PostgreSQL, executa a camada estatística determinística em um nó \texttt{function} (JavaScript), avalia em um nó \texttt{switch} se deve acionar o agente LLM via nó \texttt{http request}, valida a resposta cognitiva contra o \textit{JSON Schema}, envia os comandos de atuação para os \textit{Coils} (\texttt{FC15}) e \texttt{HR[2]} (\texttt{FC06}) através do nó \texttt{Modbus-Flex-Write} e despacha relatórios automáticos quando a severidade atingir nível de atenção ou crítico.
\end{enumerate}

\begin{figure}[htb]
  \centering
  \caption{Fluxograma completo da pipeline híbrida de aquisição, análise estatística, IA e atuação no Node-RED}
  \label{fig:fluxo-pipeline-nodered}
  \resizebox{\textwidth}{!}{%
  \begin{tikzpicture}[
    font=\sffamily\footnotesize,
    proc/.style={rectangle, rounded corners=3pt, draw=#1, fill=#1!10, thick, minimum width=3.2cm, minimum height=0.95cm, align=center, inner sep=5pt},
    dec/.style={diamond, aspect=2.1, draw=amberwarn!90!black, fill=amberwarn!12, thick, align=center, inner sep=2pt, font=\sffamily\scriptsize},
    arr/.style={-{Stealth[length=2.2mm]}, thick, draw=slateink}
  ]
    % Linha 1: Ingestão e Estatística (esquerda -> direita)
    \node[proc=slateink] (trig)   at (0,0)    {\textbf{Nó \texttt{inject} / Timer}\\Ciclo $T_p = 500\,\text{ms}$};
    \node[proc=wegblue]  (mbread) at (4.2,0)  {\textbf{\texttt{Modbus-Read}}\\Leitura \texttt{FC02} + \texttt{FC04}\\+ \texttt{FC06} Heartbeat};
    \node[proc=tealnode] (pgraw)  at (8.4,0)  {\textbf{Nó \texttt{postgresql}}\\Grava telemetria\\e consulta janela $W$};
    \node[proc=wegblue]  (stats)  at (12.6,0) {\textbf{Nó \texttt{function} (JS)}\\Calcula $\lambda_i$, $D_i$, $z_i$\\e \textit{chattering}};

    % Linha 2: Decisão, IA, Validação e Escrita (direita -> esquerda)
    \node[dec]               (decia)   at (12.6,-2.35) {Nó \texttt{switch}:\\Anomalia/sumário?};
    \node[proc=crimsonalert] (llm)     at (8.4,-2.35)  {\textbf{\texttt{http request} (LLM)}\\Diagnóstico contextual\\e recomendação};
    \node[dec]               (decjson) at (4.2,-2.35)  {JSON Schema\\válido?};
    \node[proc=wegblue]      (mbwrite) at (0,-2.35)    {\textbf{\texttt{Modbus-Flex-Write}}\\Escrita \texttt{FC15} (Coils)\\e \texttt{FC06} (\texttt{HR[2]})};

    % Linha 3: Fallback determinístico
    \node[proc=amberwarn]    (fallback) at (4.2,-4.45) {\textbf{Fallback Seguro}\\Aplica ação estatística\\determinística};

    % Linha 4: Auditoria e Relatórios Automáticos (sem cruzamentos)
    \node[proc=tealnode]     (pgdiag)   at (8.4,-6.25) {\textbf{Nó \texttt{postgresql}}\\Registra diagnóstico\\e comandos enviados};
    \node[proc=amberwarn]    (rep)      at (12.6,-6.25){\textbf{Relatório Automático}\\E-mail / Telegram\\(Markdown estruturado)};

    % Conexões Linha 1 -> Linha 2
    \draw[arr] (trig) -- (mbread);
    \draw[arr] (mbread) -- (pgraw);
    \draw[arr] (pgraw) -- (stats);
    \draw[arr] (stats) -- (decia);

    % Conexões Linha 2
    \draw[arr] (decia.west) -- node[above, font=\sffamily\scriptsize]{Sim} (llm.east);
    \draw[arr] (llm.west) -- (decjson.east);
    \draw[arr] (decjson.west) -- node[above, font=\sffamily\scriptsize]{Sim} (mbwrite.east);

    % Ramo Fallback (JSON inválido / timeout LLM)
    \draw[arr] (decjson.south) -- node[right, font=\sffamily\scriptsize]{Não} (fallback.north);
    \draw[arr] (fallback.west) -| ([xshift=0.55cm]mbwrite.south);

    % Ramo Auditoria pós-escrita Modbus (passa livre abaixo do Fallback)
    \draw[arr] ([xshift=-0.55cm]mbwrite.south) |- (pgdiag.west);

    % Ramo Nominal (sem invocação LLM) -> Auditoria
    \draw[arr] (decia.south) -- ++(0,-1.25) -| node[pos=0.22, above, font=\sffamily\scriptsize]{Não (Nominal)} (pgdiag.north);

    % Disparo de relatório automático
    \draw[arr] (pgdiag.east) -- node[above, font=\sffamily\scriptsize]{Sev. $\ge 1$} (rep.west);
  \end{tikzpicture}%
  }
  \fonte{Elaborado pelo autor (2026).}
\end{figure}

\section{Modelagem Estatística em Janela Temporal Deslizante}
\label{sec:nodered_estatistica}

A camada determinística avalia o comportamento dinâmico de cada sensor discreto $i \in \{1, 2, 3, 4\}$ sobre uma janela temporal deslizante de duração $W = 60\,\text{s}$ (equivalente a $K = W / T_p = 120$ amostras nominais). Três indicadores quantitativos são extraídos continuamente \cite{chandola_2009}:

\begin{enumerate}
  \item \textbf{Taxa de acionamento ($\lambda_i$):} obtida pela diferença modular do contador acumulado no CLP (\textit{Input Registers} \texttt{0..3}) entre o início e o término da janela, convertida em eventos por minuto:
  \begin{equation}
    \label{eq:taxa_acionamento}
    \lambda_i(W) = \frac{N_i(W)}{W} \times 60
  \end{equation}
  em que $N_i(W) = \left(C_i(t) - C_i(t - W)\right) \bmod 65536$ representa o total de bordas de subida observadas na janela.

  \item \textbf{Proporção de tempo ativo (\textit{duty cycle}, $D_i$):} quantifica a fração de amostras em que o sensor permaneceu em estado de alarme ou acionado:
  \begin{equation}
    \label{eq:duty_cycle}
    D_i(W) = \frac{1}{K} \sum_{k=1}^{K} s_{i,k}
  \end{equation}
  onde $s_{i,k} \in \{0, 1\}$ é o estado booleano filtrado do sensor $i$ na amostra $k$.

  \item \textbf{Escore padronizado de anomalia (\textit{Z-score}, $z_i$):} mede o afastamento da taxa atual de acionamentos em relação à linha de base histórica ($\mu_i, \sigma_i$) calculada sobre as últimas 24 horas de operação em regime nominal:
  \begin{equation}
    \label{eq:zscore}
    z_i(W) = \frac{\lambda_i(W) - \mu_i}{\sigma_i + \epsilon}
  \end{equation}
  sendo $\epsilon = 0{,}01\,\text{eventos/min}$ uma constante positiva de regularização que previne divisões por zero quando $\sigma_i = 0$.
\end{enumerate}

O critério estatístico classifica cada janela em quatro níveis de severidade padronizados, apresentados no Quadro~\ref{quadro:severidade}. Essa classificação determina simultaneamente o roteamento no nó \texttt{switch} para invocação do agente LLM e a ação de \textit{fallback} caso a camada cognitiva esteja indisponível.

\begin{quadro}[htb]
  \centering
  \SingleSpacing
  \caption{Matriz determinística de classificação de severidade e política de acionamento da camada híbrida}
  \label{quadro:severidade}
  \footnotesize
  \begin{tabularx}{\textwidth}{|p{1.9cm}|p{1.1cm}|L|p{3.9cm}|}
    \hline
    \textbf{Severidade} & \textbf{Código} & \textbf{Condição Estatística / Evento na Janela $W$} & \textbf{Ação no Node-RED e Atuadores} \\
    \hline
    \texttt{NORMAL} & \texttt{0} & $|z_1| < 2{,}0$; $D_2 = 0$; $D_3 = 0$; \texttt{DI3=1}; sem watchdog & Persistência regular; \texttt{L1=ON}, \texttt{L2=OFF}, \texttt{L3=OFF}; LLM apenas em sumário periódico \\
    \hline
    \texttt{ATENCAO} & \texttt{1} & $2{,}0 \le |z_1| < 3{,}0$ (desvio de cadência) ou $0 < D_2 < 0{,}30$ (nível alto intermitente) & Invoca agente LLM para diagnóstico; avalia acionamento preventivo de \texttt{L2}; registra alerta \\
    \hline
    \texttt{CRITICO} & \texttt{2} & $|z_1| \ge 3{,}0$ (\textit{chattering}/bloqueio), $D_2 \ge 0{,}30$, \texttt{DI2=1} (sobretemp.) ou watchdog ativo & Invoca agente LLM; comanda \texttt{L1=OFF}, \texttt{L2=ON}, \texttt{L3=ON}; dispara relatório imediato \\
    \hline
    \texttt{EMERGENCIA} & \texttt{3} & \texttt{DI3=0} ou \texttt{xFlag\_EmergTravada=1} & CLP já impõe veto local imediato; agente LLM gera relatório forense de incidente \\
    \hline
  \end{tabularx}
  \fonte{Elaborado pelo autor (2026).}
\end{quadro}

O Quadro~\ref{quadro:codigo_js} apresenta o código JavaScript de referência implementado no nó \texttt{function} do Node-RED para processar as linhas retornadas pelo PostgreSQL em \texttt{msg.payload}, ler a linha de base do contexto \texttt{flow}, computar as Equações~\eqref{eq:taxa_acionamento}--\eqref{eq:zscore} e preparar a mensagem para o nó \texttt{switch}.

\begin{quadro}[htb]
  \centering
  \SingleSpacing
  \caption{Código JavaScript do nó \texttt{function} no Node-RED para agregação temporal e detecção estatística}
  \label{quadro:codigo_js}
  \begin{lstlisting}[language=JavaScript]
// Entrada: msg.payload contendo array de amostras da janela W=60s vindo do no postgresql
const amostras = Array.isArray(msg.payload) ? msg.payload : [];
if (amostras.length === 0) {
  msg.payload = { status: "SEM_DADOS", acionar_llm: false };
  return msg;
}

const inicio = amostras[0];
const fim    = amostras[amostras.length - 1];
const K      = amostras.length;
const W_SEG  = 60;
const EPS    = 0.01;

// Diferenca modular segura para contadores Modbus UINT (16 bits)
const deltaMod16 = (atual, anterior) => (atual - anterior + 65536) % 65536;

const nPresenca  = deltaMod16(fim.cnt_presenca, inicio.cnt_presenca);
const taxaPres   = (nPresenca / W_SEG) * 60;
const dutyNivel  = amostras.reduce((acc, s) => acc + (s.di_nivel_alto ? 1 : 0), 0) / K;
const dutyTemp   = amostras.reduce((acc, s) => acc + (s.di_sobretemp  ? 1 : 0), 0) / K;
const emergAtiva = fim.flag_emerg_travada || !fim.di_emergencia_ok;

// Parametros de baseline lidos do contexto de fluxo (flow) ou da consulta SQL
const baseCtx   = flow.get("baseline_presenca") || { mu: 12.0, sigma: 2.5 };
const muPres    = Number(fim.baseline_mu_presenca ?? baseCtx.mu);
const sigmaPres = Number(fim.baseline_sigma_presenca ?? baseCtx.sigma);
const zPresenca = (taxaPres - muPres) / (sigmaPres + EPS);

// Classificacao deterministica de severidade e acao de fallback
let severidade = 0; // 0=NORMAL, 1=ATENCAO, 2=CRITICO, 3=EMERGENCIA
let fallbackCmd = { cmd_motor: true, cmd_valvula: false, cmd_alarme: false };

if (emergAtiva) {
  severidade = 3;
  fallbackCmd = { cmd_motor: false, cmd_valvula: false, cmd_alarme: true };
} else if (dutyTemp > 0 || dutyNivel >= 0.30 || Math.abs(zPresenca) >= 3.0 || fim.flag_watchdog) {
  severidade = 2;
  fallbackCmd = { cmd_motor: false, cmd_valvula: true, cmd_alarme: true };
} else if (dutyNivel > 0 || Math.abs(zPresenca) >= 2.0) {
  severidade = 1;
  fallbackCmd = { cmd_motor: true, cmd_valvula: true, cmd_alarme: false };
}

msg.payload = {
  janela_inicio: inicio.capturado_em,
  janela_fim: fim.capturado_em,
  amostras_k: K,
  taxa_presenca_min: Number(taxaPres.toFixed(2)),
  duty_nivel: Number(dutyNivel.toFixed(3)),
  duty_sobretemp: Number(dutyTemp.toFixed(3)),
  zscore_presenca: Number(zPresenca.toFixed(2)),
  estado_clp: fim.estado_clp,
  cod_intertravamento: fim.cod_intertrav,
  severidade_estatistica: severidade,
  fallback_comandos: fallbackCmd,
  acionar_llm: severidade >= 1 || flow.get("ciclo_sumario_horario") === true
};
return msg;
  \end{lstlisting}
  \fonte{Elaborado pelo autor (2026).}
\end{quadro}

\section{Camada Cognitiva (Agente LLM), Governança por JSON Schema e Relatórios}
\label{sec:nodered_ia}

Quando \texttt{msg.payload.acionar\_llm} assume valor verdadeiro, o nó \texttt{switch} direciona a mensagem para o ramo cognitivo, onde um nó \texttt{function} monta o \textit{prompt} estruturado contendo o resumo estatístico da janela, o estado atual do PLC500 e os últimos cinco eventos do PostgreSQL, enviando-o via nó \textbf{\texttt{http request}} à API do modelo de linguagem \cite{nodered_docs}. Essa arquitetura suporta tanto provedores em nuvem quanto servidores de inferência locais na rede interna, impondo três salvaguardas obrigatórias:

\begin{enumerate}
  \item \textbf{Restrição estrutural por \textit{JSON Schema}:} o nó validador JSON do Node-RED confronta a resposta do modelo com o esquema especificado no Quadro~\ref{quadro:json_schema} (\texttt{additionalProperties: false}, tipos booleanos estritos para os comandos de LEDs e enumeração fechada de severidade).
  \item \textbf{Regra de monotonicidade de segurança:} o nó \texttt{function} pós-validação impede que o modelo rebaixe uma severidade estatística \texttt{CRITICO} ou \texttt{EMERGENCIA} para \texttt{NORMAL} e bloqueia qualquer comando \texttt{cmd\_motor = true} enquanto houver sobretemperatura, nível alto sustentado ou emergência travada, formatando em seguida o payload para o nó \texttt{Modbus-Flex-Write} (\texttt{fc: 15, address: 0, quantity: 3}).
  \item \textbf{Relatório técnico automatizado:} nos eventos com severidade $\ge 1$ ou nos fechamentos periódicos de turno, os campos \texttt{diagnostico}, \texttt{causa\_provavel} e \texttt{resumo\_executivo} são formatados em Markdown por um nó \texttt{template} e despachados automaticamente por e-mail ou Telegram aos operadores.
\end{enumerate}

\begin{quadro}[htb]
  \centering
  \SingleSpacing
  \caption{\textit{JSON Schema} obrigatório de validação da saída do agente LLM no Node-RED}
  \label{quadro:json_schema}
  \begin{lstlisting}[language=JSON]
{
  "$schema": "https://json-schema.org/draft/2020-12/schema",
  "title": "DiagnosticoAtuacaoPLC500",
  "type": "object",
  "additionalProperties": false,
  "required": [
    "severidade", "diagnostico", "causa_provavel",
    "comandos_atuadores", "requer_relatorio", "resumo_executivo"
  ],
  "properties": {
    "severidade": {
      "type": "string",
      "enum": ["NORMAL", "ATENCAO", "CRITICO", "EMERGENCIA"]
    },
    "diagnostico": { "type": "string", "minLength": 10, "maxLength": 400 },
    "causa_provavel": { "type": "string", "minLength": 5, "maxLength": 250 },
    "comandos_atuadores": {
      "type": "object",
      "additionalProperties": false,
      "required": ["cmd_motor", "cmd_valvula", "cmd_alarme"],
      "properties": {
        "cmd_motor":   { "type": "boolean" },
        "cmd_valvula": { "type": "boolean" },
        "cmd_alarme":  { "type": "boolean" }
      }
    },
    "requer_relatorio": { "type": "boolean" },
    "resumo_executivo": { "type": "string", "minLength": 20, "maxLength": 800 }
  }
}
  \end{lstlisting}
  \fonte{Elaborado pelo autor (2026).}
\end{quadro}

\section{Persistência e Rastreabilidade em PostgreSQL}
\label{sec:nodered_postgres}

Para assegurar rastreabilidade forense completa das leituras e decisões automatizadas, utiliza-se o banco de dados relacional \textbf{PostgreSQL} acessado a partir do Node-RED \cite{postgresql_docs}. O esquema relacional separa telemetria bruta de alta cadência, métricas agregadas por janela, diagnósticos emitidos pelo agente LLM (registrando explicitamente quando o \textit{fallback} estatístico foi utilizado) e trilha de auditoria dos comandos escritos via \texttt{Modbus-Flex-Write}. O script DDL completo é apresentado no Quadro~\ref{quadro:ddl_sql}.

\begin{quadro}[htb]
  \centering
  \SingleSpacing
  \caption{Script DDL SQL (PostgreSQL) da base de telemetria, análise híbrida e auditoria de atuação}
  \label{quadro:ddl_sql}
  \begin{lstlisting}[language=PGSQL]
CREATE TABLE IF NOT EXISTS telemetria_sensores (
  id                  BIGSERIAL PRIMARY KEY,
  capturado_em        TIMESTAMPTZ NOT NULL DEFAULT NOW(),
  di_presenca         BOOLEAN NOT NULL,
  di_nivel_alto       BOOLEAN NOT NULL,
  di_sobretemp        BOOLEAN NOT NULL,
  di_emergencia_ok    BOOLEAN NOT NULL,
  flag_emerg_travada  BOOLEAN NOT NULL,
  flag_watchdog       BOOLEAN NOT NULL,
  cnt_presenca        INTEGER NOT NULL,
  cnt_nivel_alto      INTEGER NOT NULL,
  cnt_sobretemp       INTEGER NOT NULL,
  estado_clp          SMALLINT NOT NULL,
  mascara_do_real     SMALLINT NOT NULL,
  cod_intertrav       SMALLINT NOT NULL
);
CREATE INDEX IF NOT EXISTS idx_tel_capturado ON telemetria_sensores (capturado_em DESC);

CREATE TABLE IF NOT EXISTS janelas_analiticas (
  id                     BIGSERIAL PRIMARY KEY,
  janela_inicio          TIMESTAMPTZ NOT NULL,
  janela_fim             TIMESTAMPTZ NOT NULL,
  amostras_k             INTEGER NOT NULL,
  taxa_presenca_min      DOUBLE PRECISION NOT NULL,
  duty_nivel             DOUBLE PRECISION NOT NULL,
  duty_sobretemp         DOUBLE PRECISION NOT NULL,
  zscore_presenca        DOUBLE PRECISION NOT NULL,
  severidade_estatistica SMALLINT NOT NULL CHECK (severidade_estatistica BETWEEN 0 AND 3)
);

CREATE TABLE IF NOT EXISTS diagnosticos_ia (
  id                 BIGSERIAL PRIMARY KEY,
  janela_id          BIGINT NOT NULL REFERENCES janelas_analiticas (id),
  criado_em          TIMESTAMPTZ NOT NULL DEFAULT NOW(),
  severidade_final   VARCHAR(16) NOT NULL,
  origem_decisao     VARCHAR(16) NOT NULL CHECK (origem_decisao IN ('LLM', 'FALLBACK_STATS')),
  diagnostico        TEXT NOT NULL,
  causa_provavel     TEXT NOT NULL,
  payload_validado   JSONB NOT NULL,
  relatorio_enviado  BOOLEAN NOT NULL DEFAULT FALSE
);

CREATE TABLE IF NOT EXISTS comandos_atuadores (
  id                 BIGSERIAL PRIMARY KEY,
  diagnostico_id     BIGINT REFERENCES diagnosticos_ia (id),
  enviado_em         TIMESTAMPTZ NOT NULL DEFAULT NOW(),
  cmd_motor          BOOLEAN NOT NULL,
  cmd_valvula        BOOLEAN NOT NULL,
  cmd_alarme         BOOLEAN NOT NULL,
  severidade_hr2     SMALLINT NOT NULL,
  mascara_do_confirm SMALLINT,
  veto_local_clp     BOOLEAN NOT NULL DEFAULT FALSE
);
  \end{lstlisting}
  \fonte{Elaborado pelo autor (2026).}
\end{quadro}

Com o índice decrescente em \texttt{capturado\_em}, as consultas de janela deslizante de 60 segundos executadas pelos nós do Node-RED mantêm tempo de resposta baixo e previsível, facilitando também a conexão de dashboards operacionais.
